\section{引言}

直接偏好优化(DPO)已迅速成为将大语言模型对齐到人类偏好的主导范式 \cite{rafailov2023dpo}。通过免去独立奖励模型与在线采样的需求,DPO 将对齐简化为在偏好对上的一个分类式损失,使从业者能够以极低的基础设施开销微调模型。这种简洁性推动了广泛采用:截至 2025 年初,公共仓库上大多数开放权重的指令微调模型都将 DPO 或其变体用作后训练步骤 \cite{pal2024smaug}。然而,这种普遍性可能掩盖了一个潜在的脆弱性。DPO 的理论形式隐式地促使策略将概率质量集中到一小撮被偏好的补全上,这一现象被称为模式坍缩 \cite{razin2024vanishing}。尽管已有工作将该坍缩刻画为一种性能局限——产生重复或退化的输出——我们认为它构成了一个远为严重的问题:一项可被利用的安全漏洞。

DPO 下的模式坍缩已有充分记载。随着训练推进,模型的输出分布在不断收缩的模式集合周围变得尖锐,从而降低其生成过程的有效熵 \cite{razin2024vanishing, pal2024smaug}。标准做法通过依据下游基准性能校准的早停来缓解这一点,在有用性指标达到峰值时停止训练。然而,这一停止准则对那些主导对抗鲁棒性的分布性质是无感的。我们观察到,随着输出熵下降,模型行为变得越来越确定,因而对攻击者越来越可预测。一个在给定前缀条件下输出熵很低的模型,只允许一小撮可枚举的可能续写——而这恰恰是使定向利用成为可能的条件。这种对齐诱导的坍缩与对抗可预测性之间的联系,尚未在安全文献中被审视。

与此同时,关于大模型投毒的工作已证明偏好数据是一个强有力的攻击面。Pathmanathan 等人 \cite{pathmanathan2025dpo_poison} 表明,仅污染 0.5\% 的 DPO 训练对就足以注入定向的不当行为;Wang 等人 \cite{wang2024rlhfpoison} 演示了通过偏好投毒对奖励模型进行操纵。影子对齐攻击仅用 100 个对抗样本即可实现安全性退化 \cite{yang2023shadow}。另一方面,诸如 GCG \cite{zou2023gcg} 这样的越狱方法利用可访问梯度的模型来寻找诱发有害输出的对抗后缀。尚未被探索的是这两条研究脉络的交集:DPO 内在的模式坍缩如何放大投毒效力,以及由此产生的确定性如何催生一类新攻击——而这类攻击是投毒文献和越狱文献各自孤立地都未曾处理的。

本文通过一个新指标——\textit{偏好坍缩可利用度}(PCE)——对这一交集进行形式化,该指标量化了 DPO 训练中熵降低所引发的安全风险。我们的贡献有五点。第一,我们定义了 PCE,并建立其与确定性输出体制下对抗成功概率之间的理论关系。第二,我们经验性地证明 PCE 随 DPO 训练步数单调上升,并识别出一个先于标准性能峰值出现的临界转折点——这意味着模型通常会被训练到越过安全阈值之后。第三,我们设计了\textit{坍缩加速器}攻击:一种仅需 1\%--2\% 训练数据的数据投毒策略,它主动将模式坍缩引向攻击者指定的输出模式,实现超过 85\% 的定向有害生成率。第四,我们审计了六个被广泛部署的开源 DPO 模型,发现它们的 PCE 得分全部高于 0.5,处于高风险区间。第五,我们提出了一个熵正则化的 DPO 目标,在标准对齐基准上退化不足 2\% 的同时将 PCE 降低 59\%,提供了一个与现有训练流程兼容的实用防御。

本文余下部分组织如下。第~\ref{sec:background} 节回顾 DPO 的训练动力学与偏好学习的威胁模型。第~\ref{sec:pce} 节形式化 PCE 指标及其理论性质。第~\ref{sec:collapse_attack} 节给出坍缩加速器攻击的方法。第~\ref{sec:experiments} 节报告包括开源模型审计在内的实证结果。第~\ref{sec:defense} 节引入并评估我们的熵正则化防御。第~\ref{sec:discussion} 节讨论局限性与更广泛的影响,第~\ref{sec:conclusion} 节作出总结。
